\chapter{CONCLUSIONS AND FUTURE WORK}\label{ch:conclusion}

\begin{doublespacing}

\section{Summary of Findings}

This thesis has presented a unified framework for analyzing multi-robot
search and rescue across four algorithmic models that vary energy
capacity and communication assumptions. Theoretical upper bounds on
the expected finder competitive ratio were derived for each model and
validated by Monte Carlo simulation over 500 trials per configuration.

Four findings emerged, each corresponding to one of the research
hypotheses stated in Chapter~III.

\begin{enumerate}
\item Auction-based coordination yields a $6.6\times$ improvement in
finder competitive ratio over random uncoordinated search, confirming
that the communication mode---not the fleet size---is the dominant
determinant of search efficiency, consistent with the theoretical
prediction of Chrobak et al.~\cite{Chrobak2015}.
\item Probability-aware distributed auctions outperform a centralized
distance-optimal Hungarian baseline by $5.7$ percent. An ablation in
which the Hungarian solver is given the same probability-weighted
costs as the auction eliminates the gap, demonstrating that the
benefit is attributable to the bid function rather than to
centralization.
\item Multi-node sorties recover $92$ percent of the gap between the
energy-constrained single-sortie algorithm and the
unconstrained-energy auction baseline. A single greedy chain extension
per sortie visits $\sim\!5.3$ sites on average at the default setting,
cutting the iteration count from $4.81$ to $1.38$.
\item A bid function $p_i / d^2$ that penalizes distance quadratically
improves on the standard $p_i / d$ bid by $4.2$ percent under energy
constraints. The result is justified by a cost-foreclosure argument:
distant anchor sites not only cost more energy but also shorten the
chain of future extensions, so the effective cost of distance is
superlinear.
\end{enumerate}

All four theoretical bounds held empirically across every trial. The
bounds on Models~1 and~2 are relatively tight; the Model~4 bound is
loose because the analytical estimate of sites per sortie
underestimates what the greedy chain actually achieves.

\section{Future Work}

Several directions for extending this work are natural.

\textbf{Alternative objectives.} Finder competitive ratio measures
only the finding robot's distance; it says nothing about the other
robots' energy expenditure. Minimax workload balancing, total-fleet
distance, or time-to-last-cleared-site are alternative objectives that
would favor different bid functions. A systematic study of the
trade-off among these objectives on a common benchmark is an open
problem.

\textbf{Tighter Model 4 bound.} The looseness of the Model 4 bound
comes from the pessimistic estimate $S = 1 + \lfloor (E -
2\,d_{\mathrm{avg}}) / d_{\mathrm{hop}} \rfloor$. A more refined
analysis would compute the expected length of a greedy chain under the
prior distribution, not the worst-case length under a uniform
distribution, and could plausibly produce a bound within a factor of
two of the empirical mean.

\textbf{Fleet-size scaling conjecture.} The robot sweep suggests that
coordinated models achieve FCR $= O(1)$ as $R \to \infty$ while
uncoordinated models achieve only FCR $= \Omega(R^{-1})$. A formal
proof of this conjecture would strengthen the case for coordination
as the primary design lever.

\textbf{Relaxing assumptions.} Perfect sensors, flat terrain, and
uniform energy budgets are idealizations. Detection probability less
than one, shortest-path distance over an obstacle graph, and per-robot
energy heterogeneity are realistic generalizations. The auction
mechanism adapts to each of these without modification; what changes
is the structure of the bounds.

\textbf{Physical deployment.} The simulation validates the algorithms
under the stated assumptions, but physical deployment introduces
additional concerns: sensor noise, communication loss, obstacle
avoidance, and motor control variability. Porting the algorithms to a
small team of UAVs or ground robots and measuring FCR on physical
instances would provide the strongest test of the framework and is
the logical next step.

\end{doublespacing}
